\chapter{Measles}

\section*{Definition and Overview}
Measles (rubeola) is a highly contagious, acute viral respiratory infection caused by the measles virus, a single-stranded, negative-sense RNA virus of the genus \textit{Morbillivirus} within the family \textit{Paramyxoviridae}. It is characterised by a prodromal phase of high fever, coryza, cough, conjunctivitis, and pathognomonic Koplik's spots, followed by a generalised maculopapular rash that spreads cephalocaudally. Measles is one of the most infectious diseases known to medicine, transmitted via respiratory droplets or airborne aerosols that can remain infective in the air for up to two hours. Despite the availability of a highly effective vaccine (often combined as the measles, mumps, and rubella vaccine), measles remains a major cause of morbidity and mortality globally, particularly in young children and in areas with low vaccination coverage.

\section*{Epidemiology}
Globally, measles remains a leading vaccine-preventable cause of death in infants and young children, particularly in low-income nations with suboptimal health infrastructure. In Australia, endemic transmission of measles was declared eliminated in 2014 due to high herd immunity achieved through the National Immunisation Program. However, sporadic cases and localised outbreaks continue to occur, imported by unvaccinated travellers returning from endemic areas. Measles has a basic reproduction number estimated between 12 and 18, meaning a single infected individual can transmit the virus to 12--18 susceptible people. The risk of severe complications and death is highest in children under 5 years, adults over 20 years, pregnant women, and immunocompromised individuals.

\section*{Aetiology and Pathophysiology}
Measles is caused by the measles virus. The pathophysiology involves primary replication in the respiratory epithelium followed by systemic dissemination.
\begin{itemize}
    \item \textbf{Viral entry and replication:} The measles virus binds to host cellular receptors (specifically CD150/SLAM on immune cells and nectin-4 on epithelial cells) in the respiratory tract. It replicates locally and in regional lymph nodes.
    \item \textbf{Primary and secondary viraemia:} Within days, the virus enters the bloodstream (primary viraemia), spreading to the reticuloendothelial system. A subsequent secondary viraemia distributes the virus to the skin, respiratory tract, viscera, and central nervous system.
    \item \textbf{Immunosuppression:} The virus infects and destroys CD4+ and CD8+ T-lymphocytes, B-lymphocytes, and dendritic cells, causing a profound, transient state of immune suppression (immune amnesia) that lasts for months to years, leaving the host highly susceptible to secondary bacterial infections.
    \item \textbf{Pathognomonic lesions:} The rash is a result of a delayed-type hypersensitivity reaction (cellular immune response) to the viral antigen deposited in the skin capillaries.
\end{itemize}

\section*{Clinical Features}
The clinical course of measles progresses through distinct phases: incubation, prodrome, exanthem, and recovery.
\begin{itemize}
    \item An incubation period of 10 to 14 days, during which the patient is asymptomatic.
    \item A prodromal phase lasting 2 to 4 days, characterised by high fever (often reaching >40 degrees Celsius), cough, coryza (runny nose), and conjunctivitis (the classic "three Cs").
    \item Koplik's spots, which are pathognomonic, small, bluish-white spots on a red background on the buccal mucosa opposite the lower molars, appearing 1 to 2 days before the rash.
    \item A maculopapular rash appearing 3 to 5 days after the onset of prodromal symptoms, beginning on the face and behind the ears, then spreading down the neck, trunk, and limbs.
    \item The rash initially blanches on pressure, becomes confluent on the face and trunk, and gradually fades in the order of appearance, often leaving a brown discolouration and fine desquamation.
    \item Constitutional symptoms like anorexia, generalised malaise, diarrhoea (especially in children), and cervical lymphadenopathy.
\end{itemize}

\section*{Differential Diagnosis}
The clinical features of measles can overlap with other childhood exanthems and viral infections:
\begin{itemize}
    \item \textbf{Rubella:} Causes a similar maculopapular rash and fever, but rubella is characterised by milder symptoms, prominent post-auricular and suboccipital lymphadenopathy, and lacks Koplik's spots.
    \item \textbf{Roseola infantum:} Presents with high fever that abruptly subsides before the onset of a maculopapular rash, which begins on the trunk and spreads to the limbs, opposite to the pattern of measles.
    \item \textbf{Scarlet fever:} Caused by Group A Streptococcus; features a diffuse, rough, sandpaper-like rash, strawberry tongue, and tonsillitis, but lacks the cough, conjunctivitis, and Koplik's spots.
    \item \textbf{Kawasaki disease:} An inflammatory vasculitis presenting with prolonged fever, conjunctivitis, rash, and mucosal changes, but distinguished by strawberry tongue, swollen hands/feet, and cardiac investigations.
    \item \textbf{Drug eruption:} An allergic reaction to medications that can cause a similar maculopapular rash, but lacks the viral prodrome, Koplik's spots, and infectious symptoms.
\end{itemize}

\section*{When to Seek Medical Care}
If a parent suspects their child has measles, they should contact their doctor's clinic before visiting to prevent exposure to other patients in the waiting room.

\textbf{Call 000 (or local emergency services) for:}
\begin{itemize}
    \item Difficulty breathing, rapid breathing, or grunting (suggesting severe pneumonia).
    \item Severe lethargy, confusion, stiff neck, or seizures (suggesting encephalitis).
    \item A high fever accompanied by severe, uncontrollable vomiting or signs of dehydration.
    \item Severe bleeding or bruising.
\end{itemize}

Other non-emergency reasons to contact a doctor include:
\begin{itemize}
    \item Suspected exposure to someone with measles, especially if the person is unvaccinated, pregnant, or immunocompromised.
    \item High fever that does not settle with paracetamol or ibuprofen.
    \item Worsening ear pain, irritability, or fluid draining from the ear.
    \item Persistent cough or chest congestion.
\end{itemize}

\section*{Diagnosis and Investigations}
Diagnosis is often suspected clinically but must be confirmed microbiologically for public health reporting.
\begin{itemize}
    \item \textbf{Measles serology:} Detection of measles-specific IgM antibodies in the blood, which indicates acute or recent infection. IgG antibodies indicate past infection or vaccination.
    \item \textbf{Reverse transcription PCR:} The most sensitive and specific test, performed on nasopharyngeal swabs, throat swabs, or urine samples, to detect measles virus RNA.
    \item \textbf{Viral culture:} Historically used but now largely replaced by PCR due to the latter's speed and accuracy.
    \item \textbf{Full blood count:} Often shows leucopenia and lymphopenia, which are typical of acute viral infections.
    \item \textbf{Chest X-ray:} Indicated if secondary bacterial or primary viral pneumonia is suspected, showing diffuse interstitial infiltrates or lobar consolidation.
\end{itemize}

\section*{Management}
There is no specific antiviral treatment for measles. Management is primarily supportive, aimed at relieving symptoms and treating complications.
\begin{itemize}
    \item \textbf{Symptomatic relief:} Adequate hydration, rest, and antipyretics (paracetamol or ibuprofen) to control fever and discomfort; aspirin must be avoided in children due to the risk of Reye's syndrome.
    \item \textbf{Vitamin A supplementation:} Recommended by the World Health Organization for all children diagnosed with acute measles. High-dose vitamin A reduces the risk of blindness, severe croup, and mortality.
    \item \textbf{Antibiotic therapy:} Indicated only for secondary bacterial infections, such as otitis media or bacterial pneumonia.
    \item \textbf{Isolation and infection control:} Patients are highly infectious from 4 days before to 4 days after the onset of the rash. Isolation in negative-pressure rooms is required in hospital settings.
    \item \textbf{Post-exposure prophylaxis:} Unvaccinated contacts can receive the MMR vaccine within 72 hours of exposure, or human normal immunoglobulin within 6 days of exposure, to prevent or attenuate the disease.
    \item \textbf{Public health notification:} Measles is a notifiable disease in Australia; clinicians must immediately report suspected cases to the local public health unit.
\end{itemize}

\section*{Complications}
\begin{itemize}
    \item \textbf{Otitis media:} The most common complication of measles, occurring in up to 10\% of cases, which can lead to permanent hearing loss.
    \item \textbf{Pneumonia:} Severe viral or secondary bacterial pneumonia, which is the leading cause of measles-related deaths in young children.
    \item \textbf{Encephalitis:} Acute inflammation of the brain occurring in approximately 1 in 1,000 cases, causing permanent neurological damage or death.
    \item \textbf{Subacute sclerosing panencephalitis:} A rare, fatal, progressive degenerative neurological disorder of the central nervous system that develops 7--10 years after measles infection.
    \item \textbf{Gastrointestinal complications:} Severe diarrhoea and vomiting, leading to dehydration and malnutrition, particularly in developing countries.
    \item \textbf{Ocular complications:} Keratitis and corneal ulceration, which can progress to permanent blindness, especially in children with vitamin A deficiency.
    \item \textbf{Pregnancy complications:} Increased risk of miscarriage, stillbirth, premature labour, and low birth weight.
\end{itemize}

\section*{Prognosis and Outcomes}
The prognosis of measles is generally excellent in well-nourished, vaccinated populations, with most patients recovering fully within 10--14 days. However, in malnourished children, particularly those with vitamin A deficiency or compromised immune systems, the disease can be severe, with global case-fatality rates reaching 1--3\% in developing areas. The development of severe complications, such as pneumonia or encephalitis, significantly worsens the prognosis. Subacute sclerosing panencephalitis is universally fatal but extremely rare.

\section*{Prevention and Self-Care}
\begin{itemize}
    \item Get vaccinated with two doses of the MMR vaccine, which is the most effective way to prevent measles. In Australia, these are scheduled at 12 and 18 months under the National Immunisation Program.
    \item Check your vaccination status if you were born after 1966 and are planning to travel overseas, as many adults require a booster dose.
    \item Isolate the infected individual at home, keeping them away from school, work, and public places until at least 4 days after the rash appears.
    \item Encourage the patient to drink plenty of water, clear broths, or oral rehydration solutions to prevent dehydration.
    \item Use a cool-mist humidifier or saline nose drops to help soothe a persistent cough and clear nasal passages.
    \item Keep the room dimly lit if the patient experiences photophobia (sensitivity to light) due to conjunctivitis.
    \item Ensure any exposed, unvaccinated household contacts consult their doctor immediately regarding post-exposure vaccination or immunoglobulin treatment.
\end{itemize}

\section*{Sources and Further Reading}
The following organisations provide authoritative, evidence-based information on measles and MMR vaccination:
\begin{itemize}
    \item Healthdirect Australia. \textit{Measles symptoms, diagnosis, and MMR vaccination.} https://www.healthdirect.gov.au/measles
    \item World Health Organization. \textit{Measles fact sheet and global immunisation data.} https://www.who.int/news-room/fact-sheets/detail/measles
    \item Centers for Disease Control and Prevention. \textit{Measles clinical features and vaccination guidance.} https://www.cdc.gov/measles/
    \item NHS. \textit{Measles overview, prevention, and treatment.} https://www.nhs.uk/conditions/measles/
    \item Australian Government Department of Health and Aged Care. \textit{Measles national guidelines for public health units.} https://www.health.gov.au/
\end{itemize}
